\chapter{报告概述}

本报告研究六类典型地学智能任务：断层预测、地震相分割、储层物性预测、岩相分类、
甜点评价和三维重建。六个赛道的数据形态不同，评价目标也不同，因此本项目没有用
一个模型替代所有基线，而是为每个任务选择与数据结构相匹配的模型，再检验基础模型
能否提供稳定、可归因的增益。

报告的主线只有三步。第一步定义任务、输入和输出，并给出模型的数学原理。第二步说明
数据划分、参数优化和模型选择方式。第三步在固定评价指标下报告实验结果。每个赛道
均采用“任务背景、方法原理、训练策略、实验设计、评估指标、实验结果”六个固定章节，
便于读者横向比较。

实验表明，基础模型的价值取决于接入位置。SAM2 交叉注意力在地震相开发集上提高了
平均交并比，但同架构随机初始化对照尚未完成，因此增益暂不能完全归因于预训练权重。
Chronos-2 在甜点评价的 T3 开发实验中优于归档 XGBoost 和因果历史均值。TabICL
在三个物性目标的开发集上改善了 RMSE，但仍缺少同架构随机权重对照；MOMENT 的
同架构对照已经完成，预训练模型略优于随机初始化，却明显弱于 XGBoost。断层预测
受制于连续三维开发数据不足，本报告仅给出二维基线结果。

三维重建采用了不同的基础模型接入方式。OpenMind-MAE 不再直接替代克里金，也不承担
不稳定的残差回归，而是以冻结表征参与邻域度量。最终 P21 固定三核集成的开发 OOF
RMSE 为 $0.027734$，低于普通克里金的 $0.028450$；在预注册的同场区历史版本
迁移测试中，RMSE 又相对 PyKrige 降低 $1.4529\%$，五个空间折中四折改善。
这一证据支持“基础模型表征在固定地统计邻域中有用”，但不外推为跨场区或赛事隐藏
测试结论。

这些结果说明，模型规模本身不是研究终点，接入机制和对照实验才决定结果。
报告因此同时呈现改进、未改进和数据不足的结果，使每项结论都与实际实验条件对应。

本版进一步把评价从单一分数扩展到空间结构和物理一致性。断层解释被放回原始地震
剖面与 UTM--TWT 场景中核对；储层物性增加 PHIF--KLOGH 联合关系；甜点生产率
增加实测--预测散点与残差分布；三维重建增加正交切片和方向变差函数。新增图件没有
改变已有模型，而是回答一个更严格的问题：模型是否同时保留点值精度、地质结构和
跨变量关系。
